\documentclass[../../main.tex]{subfiles}
\begin{document}

{\bf [解答]}

\begin{figure}[htb]
  \centering
  \begin{tikzpicture}[x=1.8cm, y=1.5cm, >=stealth,
    every node/.style={font=\small}]
    \coordinate (O) at (0, 0);
    \coordinate (C) at (1, 0);
    \coordinate (A) at (3, 0);
    \coordinate (P) at (3, 1.5);
    \coordinate (Fa) at (0, 1.5);

    % 座標軸
    \draw[->, thick, gray!80!black] (-0.3, 0) -- (4.5, 0) node[right] {$x$};
    \draw[->, thick, gray!80!black] (0, -0.4) -- (0, 2.6) node[above] {$y$};
    \node[below left=2pt] at (O) {$O$};

    % 補助線
    \draw[dashed, gray!60] (Fa) -- (P) -- (A);

    % 接線
    \draw[line width=1.0pt, teal!75!black] (0.5, -0.375) -- (4.2, 2.4)
    node[below right, font=\footnotesize] {接線};

    % 曲線 y = f(x)
    \draw[line width=1.2pt, blue!85!black, domain=0.8:4.0, samples=80, smooth]
    plot (\x, {0.2*(\x-3)*(\x-3) + 0.75*\x - 0.75})
    node[above left] {$y=f(x)$};

    % 点
    \node[circle, fill=teal!75!black, inner sep=1.4pt] at (C) {};
    \node[circle, fill=red!75!black, inner sep=1.4pt] at (P) {};

    % ラベル
    \node[below=2pt] at (C) {$c$};
    \node[below=2pt] at (A) {$a$};
    \node[left=2pt] at (Fa) {$f(a)$};
    \node[above left=2pt] at (P) {$P(a, f(a))$};
  \end{tikzpicture}
  \caption{曲線 $y=f(x)$ と点 $P(a, f(a))$ における接線}
  \label{fig:1970_1_tangent}
\end{figure}

$y=f(x)$の$x=a$における接線は
\begin{align}
  y = f'(a)(x-a) + f(a) \label{eq:1}
\end{align}
だから，これが$(c,0)$を通ることにより，
\begin{align}
  0 = f'(a)(c-a) + f(a) \label{eq:2}
\end{align}
である．

また $g(x)$ の定義から一階微分は
\begin{align}
  g'(x) = \dfrac{f'(x)(x-c) - f(x)}{(x-c)^2}
\end{align}
だから，$x=a$を代入して
\begin{align}
  g'(a) = \dfrac{f'(a)(a-c) - f(a)}{(a-c)^2} = 0 \quad (\because \text{\cref{eq:2}})
\end{align}
である．

\newpage

\end{document}
